\section{Problem 50: derivatives of the totient distribution function (Round 9)}

\subsection{1) ROUND-9 OBJECTIVE}

\textbf{Path (C) (obstruction / correction).}
Round~8 produced a dense $G_\delta$ set
\[
\Omega\subset (0,1)\setminus\mathcal E,
\qquad
\mathcal E:=\Bigl\{\frac{\varphi(n)}{n}:n\ge 1\Bigr\},
\]
on which the finite one-sided secant spectra are maximal:
\[
\Sigma_-(x)=\Sigma_+(x)=[0,\infty)
\qquad(x\in\Omega).
\]
Equivalently, every positive one-sided secant slope is realized exactly, on both sides, at arbitrarily small scales around every $x\in\Omega$.

That left one important asymmetry in the obstruction picture: Round~8 was primarily \emph{topological} (dense $G_\delta$), but it did not show that the same one-sided blow-up phenomenon is also \emph{measure-theoretically full} for the natural singular measure $\mu=dG$.

The new objective is therefore to strengthen the one-sided obstruction from ``residual'' to ``residual \emph{and} $\mu$-full.''
We prove that both sets
\[
B_-:=\Bigl\{x\in(0,1):\overline D_-G(x)=+\infty\Bigr\},
\qquad
B_+:=\Bigl\{x\in(0,1):\overline D_+G(x)=+\infty\Bigr\},
\]
are not only dense $G_\delta$ subsets of $(0,1)$ (Round~7), but in fact satisfy
\[
\mu(B_-)=\mu(B_+)=1.
\]
Hence
\[
\Xi:=\bigl(B_-\cap B_+\bigr)\setminus\mathcal E
\]
is a set that is simultaneously \emph{dense $G_\delta$}, \emph{$\mu$-full}, and disjoint from the special value set $\mathcal E$, and at every point of $\Xi$ both one-sided upper Dini derivatives are infinite.

This is a genuine strengthening of Round~8: any hypothetical point where $G'(x)$ exists and is a finite positive number must now avoid a set that is simultaneously residual and conull for the intrinsic distribution measure.

\subsection{2) Round-8 FOUNDATION USED}

We use the following earlier-round results as black boxes.

\begin{itemize}
\item \textbf{Round~1.} The limiting distribution function
\[
G(u):=\lim_{X\to\infty}\frac1X\#\Bigl\{n\le X:\frac{\varphi(n)}n\le u\Bigr\}
\qquad(0\le u\le1)
\]
exists, is continuous and strictly increasing on $[0,1]$, and the associated Borel probability measure
\[
\mu:=dG
\]
is purely singular with respect to Lebesgue measure $\lambda$. In particular,
\[
G'(x)=0\qquad\text{for Lebesgue-a.e. }x\in(0,1).
\]

\item \textbf{Round~3.} The value set
\[
\mathcal E:=\Bigl\{\frac{\varphi(n)}{n}:n\ge1\Bigr\}
\]
is countable and dense in $[0,1]$. Since $G$ is continuous, $\mu$ has no atoms; therefore
\[
\mu(\mathcal E)=0.
\]

\item \textbf{Round~7.} The sets
\[
B_-:=\Bigl\{x\in(0,1):\overline D_-G(x)=+\infty\Bigr\},
\qquad
B_+:=\Bigl\{x\in(0,1):\overline D_+G(x)=+\infty\Bigr\},
\]
where
\[
\overline D_-G(x):=\limsup_{h\to0^+}\frac{G(x)-G(x-h)}{h},
\qquad
\overline D_+G(x):=\limsup_{h\to0^+}\frac{G(x+h)-G(x)}{h},
\]
are dense $G_\delta$ subsets of $(0,1)$.

\item \textbf{Round~8.} There exists a dense $G_\delta$ set $\Omega\subset (0,1)\setminus\mathcal E$ such that every positive one-sided secant slope is realized exactly, on both sides, at arbitrarily small scales around each $x\in\Omega$.
\end{itemize}

\subsection{3) NEW INSIGHT / TOOL (ROUND-9)}

The new tool is a one-sided strengthening of the singular-density lemma used in Round~4.

\begin{quote}
\emph{If $\mu\perp\lambda$ is a finite Borel measure on $\mathbb R$, then for $\mu$-almost every $x$ both one-sided upper Lebesgue densities are infinite:}
\[
\limsup_{h\to0^+}\frac{\mu((x-h,x])}{h}=+\infty,
\qquad
\limsup_{h\to0^+}\frac{\mu([x,x+h))}{h}=+\infty.
\]
\end{quote}

Applied to $\mu=dG$, this yields
\[
\mu(B_-)=\mu(B_+)=1.
\]
Combining this with the Round~7 fact that $B_-$ and $B_+$ are dense $G_\delta$, we obtain a new set
\[
\Xi:=\bigl(B_-\cap B_+\bigr)\setminus\mathcal E
\]
that is simultaneously
\begin{itemize}
\item dense $G_\delta$,
\item $\mu$-full,
\item and disjoint from $\mathcal E$,
\end{itemize}
on which both one-sided upper Dini derivatives are infinite.

This ``residual + conull'' obstruction was not available in any previous round.

\subsection{4) ATTACK PLAN (ROUND-9)}

After Round~8, the exact remaining gap in the obstruction landscape was this:
one-sided upper blow-up was known to hold on a residual set, but not yet on a $\mu$-full set separately on each side.

We therefore prove three precise claims.

\begin{enumerate}
\item \textbf{One-sided singular-density lemma.}
For every finite singular measure $\mu\perp\lambda$ on $\mathbb R$,
\[
\limsup_{h\to0^+}\frac{\mu((x-h,x])}{h}=+\infty,
\qquad
\limsup_{h\to0^+}\frac{\mu([x,x+h))}{h}=+\infty
\]
for $\mu$-almost every $x$.

\item \textbf{Application to the totient distribution.}
For $\mu=dG$ this implies
\[
\mu(B_-)=\mu(B_+)=1.
\]

\item \textbf{Residual + conull obstruction.}
Since Round~7 already gives that $B_-$ and $B_+$ are dense $G_\delta$, the set
\[
\Xi:=\bigl(B_-\cap B_+\bigr)\setminus\mathcal E
\]
is dense $G_\delta$ and $\mu$-full. Any point where $G'(x)$ exists as a finite real number must lie outside $\Xi$; hence the candidate set is contained in a set that is simultaneously meager and $\mu$-null (and also Lebesgue-null by Round~1).
\end{enumerate}

This overcomes the Round~8 limitation by making the one-sided blow-up obstruction both topologically generic and measure-theoretically full.

\subsection{5) WORK (ROUND-9)}

\subsubsection*{5.1. A one-sided density lemma for singular measures}

Let $\mu$ be a finite Borel measure on $\mathbb R$. Define the one-sided upper densities
\[
\Theta_-^*(\mu,x):=\limsup_{h\to0^+}\frac{\mu((x-h,x])}{h},
\qquad
\Theta_+^*(\mu,x):=\limsup_{h\to0^+}\frac{\mu([x,x+h))}{h}.
\]

\begin{lemma}[One-sided upper densities of a singular measure]\label{lem:onesided-density}
If $\mu\perp\lambda$, then
\[
\Theta_-^*(\mu,x)=\Theta_+^*(\mu,x)=+\infty
\qquad\text{for }\mu\text{-almost every }x\in\mathbb R.
\]
Equivalently, for every $M>0$,
\[
\mu\bigl(\{x:\Theta_-^*(\mu,x)\le M\}\bigr)=0,
\qquad
\mu\bigl(\{x:\Theta_+^*(\mu,x)\le M\}\bigr)=0.
\]
\end{lemma}

\begin{proof}
We prove the left-sided statement; the right-sided proof is identical.

Fix $M>0$ and set
\[
A_M^-:=\{x\in\mathbb R:\Theta_-^*(\mu,x)\le M\}.
\]
Because $\mu\perp\lambda$, there exists a Borel set $S\subset\mathbb R$ such that
\[
\lambda(S)=0,
\qquad
\mu(\mathbb R\setminus S)=0.
\]
Hence $\mu(A_M^-)=\mu(A_M^-\cap S)$, so it is enough to prove
\[
\mu(A_M^-\cap S)=0.
\]

Fix $\eta>0$. Since $\lambda(S)=0$, choose an open set $O\supset S$ with
\[
\lambda(O)<\eta.
\]

For each $x\in A_M^-\cap S$, the condition $\Theta_-^*(\mu,x)\le M$ implies that there exists $r_x>0$ such that
\[
\mu((x-h,x])\le (M+1)h
\qquad(0<h<r_x).
\]
Shrinking $r_x$ if necessary, we may also ensure
\[
(x-r_x,x]\subset O.
\]
Define the interval
\[
I_x:=(x-r_x,x].
\]
Then
\begin{equation}\label{eq:Ix-left-bound}
\mu(I_x)\le (M+1)|I_x|.
\end{equation}

Now apply the one-dimensional Besicovitch covering theorem for intervals: from the family $\{I_x:x\in A_M^-\cap S\}$ one can extract a countable subfamily $\{I_j\}_{j\ge1}$ that covers $A_M^-\cap S$ and whose overlap is bounded by $2$, i.e.
no point of $\mathbb R$ belongs to more than $2$ of the intervals $I_j$.

Using \eqref{eq:Ix-left-bound} and the overlap bound, we obtain
\[
\mu(A_M^-\cap S)
\le \sum_{j\ge1}\mu(I_j)
\le (M+1)\sum_{j\ge1}|I_j|
\le 2(M+1)\,\lambda\Bigl(\bigcup_{j\ge1}I_j\Bigr).
\]
Since each $I_j\subset O$, we have
\[
\lambda\Bigl(\bigcup_{j\ge1}I_j\Bigr)\le \lambda(O)<\eta,
\]
so
\[
\mu(A_M^-\cap S)\le 2(M+1)\eta.
\]
Letting $\eta\to0$ yields $\mu(A_M^-\cap S)=0$, hence $\mu(A_M^-)=0$.

This proves that $\Theta_-^*(\mu,x)=+\infty$ for $\mu$-almost every $x$.
The proof for $\Theta_+^*(\mu,x)$ is identical, using intervals of the form $[x,x+r_x)$.
\end{proof}

\subsubsection*{5.2. Application to the totient distribution function}

Let $G$ be the distribution function from Round~1 and let $\mu=dG$.
For $x\in(0,1)$ we have
\[
\mu((x-h,x])=G(x)-G(x-h),
\qquad
\mu([x,x+h))=G(x+h)-G(x),
\]
so the one-sided upper densities of $\mu$ are precisely the one-sided upper Dini derivatives of $G$:
\[
\Theta_-^*(\mu,x)=\overline D_-G(x),
\qquad
\Theta_+^*(\mu,x)=\overline D_+G(x).
\]

\begin{theorem}[One-sided upper Dini derivatives are infinite $\mu$-a.e.]\label{thm:mu-full-onesided}
For the totient distribution measure $\mu=dG$,
\[
\overline D_-G(x)=\overline D_+G(x)=+\infty
\qquad\text{for }\mu\text{-almost every }x\in(0,1).
\]
Equivalently,
\[
\mu(B_-)=\mu(B_+)=1.
\]
In particular, the set of points where the left derivative exists as a finite real number is $\mu$-null,
and the set of points where the right derivative exists as a finite real number is also $\mu$-null.
\end{theorem}

\begin{proof}
By Round~1, $\mu\perp\lambda$. Applying Lemma~\ref{lem:onesided-density} to $\mu$ yields
\[
\Theta_-^*(\mu,x)=\Theta_+^*(\mu,x)=+\infty
\qquad\text{for }\mu\text{-almost every }x.
\]
Using the identities
\[
\Theta_-^*(\mu,x)=\overline D_-G(x),
\qquad
\Theta_+^*(\mu,x)=\overline D_+G(x),
\]
we obtain
\[
\overline D_-G(x)=\overline D_+G(x)=+\infty
\qquad\text{for }\mu\text{-almost every }x.
\]
That is exactly the statement that $\mu(B_-)=\mu(B_+)=1$.

If the left derivative (respectively right derivative) exists finitely at $x$, then the corresponding upper Dini derivative is finite, so such points lie in $(0,1)\setminus B_-$ (respectively $(0,1)\setminus B_+$). Since these complements are $\mu$-null, the stated corollaries follow.
\end{proof}

\subsubsection*{5.3. A set that is simultaneously residual, $\mu$-full, and outside $\mathcal E$}

\begin{corollary}[A residual and $\mu$-full one-sided blow-up set off $\mathcal E$]\label{cor:Xi}
Define
\[
\Xi:=\bigl(B_-\cap B_+\bigr)\setminus\mathcal E.
\]
Then:
\begin{enumerate}
\item $\Xi$ is a dense $G_\delta$ subset of $(0,1)$;
\item $\mu(\Xi)=1$;
\item for every $x\in\Xi$ one has
\[
\overline D_-G(x)=\overline D_+G(x)=+\infty.
\]
\end{enumerate}
Thus there exists a set off the special value set $\mathcal E$ that is simultaneously topologically generic and measure-theoretically full, and on which both one-sided upper Dini derivatives blow up.
\end{corollary}

\begin{proof}
By Round~7, both $B_-$ and $B_+$ are dense $G_\delta$ subsets of $(0,1)$.
Hence
\[
B_-\cap B_+
\]
is dense $G_\delta$.
Since $\mathcal E$ is countable, its complement $(0,1)\setminus\mathcal E$ is also dense $G_\delta$.
Therefore
\[
\Xi=(B_-\cap B_+)\cap\bigl((0,1)\setminus\mathcal E\bigr)
\]
is dense $G_\delta$.

By Theorem~\ref{thm:mu-full-onesided},
\[
\mu(B_-)=\mu(B_+)=1.
\]
Hence
\[
\mu(B_-\cap B_+)=1.
\]
By Round~3 and continuity of $G$, $\mu(\mathcal E)=0$. Therefore
\[
\mu(\Xi)=\mu(B_-\cap B_+)-\mu\bigl((B_-\cap B_+)\cap\mathcal E\bigr)=1.
\]
The last claim is immediate from the definition of $\Xi$.
\end{proof}

\subsubsection*{5.4. Sharpened reduction of Erd\H{o}s's question}

Define the candidate set for a finite positive derivative by
\[
D_+:=\{x\in(0,1): G'(x)\text{ exists and }0<G'(x)<\infty\}.
\]

\begin{corollary}[Any candidate point is quadruple-exceptional]\label{cor:quadruple-exceptional}
One has
\[
D_+\subseteq (0,1)\setminus\Xi.
\]
Consequently, every $x\in D_+$ satisfies all of the following:
\begin{enumerate}
\item $x\notin\mathcal E$;
\item $x\notin B_-$ and $x\notin B_+$, so both one-sided upper Dini derivatives are finite;
\item $x$ lies in a meager set;
\item $x$ lies in a $\mu$-null set;
\item $x$ lies in a Lebesgue-null set.
\end{enumerate}
In particular, the candidate set $D_+$ is contained in a set that is simultaneously meager, $\mu$-null, and Lebesgue-null.
\end{corollary}

\begin{proof}
If $x\in D_+$, then the derivative exists as a finite real number, so both one-sided upper Dini derivatives are finite.
Hence $x\notin B_-$ and $x\notin B_+$, which implies $x\notin \Xi$.
This proves
\[
D_+\subseteq (0,1)\setminus\Xi.
\]

Item (1) is already known from Round~3: every point of $\mathcal E$ has infinite left derivative.
Item (2) is exactly what we just observed.
Since $\Xi$ is dense $G_\delta$ by Corollary~\ref{cor:Xi}, its complement is meager, proving (3).
Since $\mu(\Xi)=1$, the complement is $\mu$-null, proving (4).
Finally, by Round~1, $G'(x)=0$ for Lebesgue-a.e. $x$, so any point with $0<G'(x)<\infty$ belongs to a Lebesgue-null set; this gives (5).
\end{proof}

\subsubsection*{5.5. A newly ruled-out strategy class}

Round~8 ruled out strategies based on topological regularity: on a dense $G_\delta$ set, every positive one-sided secant slope occurs exactly at arbitrarily small scales.
Round~9 adds a different obstruction: both one-sided upper Dini derivatives are infinite on a set that is simultaneously residual and $\mu$-full.

Therefore no proof strategy of the following form can succeed:
\begin{quote}
``First find a non-negligible set (either residual or of positive $\mu$-measure) on which one of the one-sided upper quotients is uniformly bounded above; then deduce the existence of a finite positive derivative from that control.''
\end{quote}
Indeed, each one-sided upper quotient already blows up on a set that is both topologically generic and conull for the intrinsic distribution measure.

\subsection{6) ADVERSARIAL VERIFICATION}

\begin{itemize}
\item \textbf{Does the one-sided density lemma really follow from singularity alone?}
Yes. The proof uses only the existence of a supporting set $S$ of Lebesgue measure $0$ and a one-dimensional Besicovitch covering argument for intervals. No arithmetic input is needed.

\item \textbf{Are half-open intervals legitimate in the covering step?}
Yes. The one-dimensional Besicovitch theorem applies to families of intervals; endpoint conventions do not affect the overlap bound or the covering argument.

\item \textbf{Did earlier rounds already imply $\mu(B_-)=\mu(B_+)=1$?}
No. Round~4 gave only the \emph{symmetric} statement
\[
\limsup_{h\to0^+}\frac{G(x+h)-G(x-h)}{2h}=+\infty
\qquad(\mu\text{-a.e. }x),
\]
and Round~7 gave only that $B_-$ and $B_+$ are dense $G_\delta$. Theorem~\ref{thm:mu-full-onesided} is a genuine one-sided measure-theoretic strengthening.

\item \textbf{Why is $\mu(\mathcal E)=0$?}
Because $\mathcal E$ is countable (Round~3) and $G$ is continuous (Round~1), so $\mu$ has no atoms.

\item \textbf{Could a point in $\Xi$ still have a finite two-sided derivative?}
No. A finite derivative would force both one-sided upper Dini derivatives to be finite, contradicting the defining property of $\Xi$.
\end{itemize}

\subsection{7) FINAL (EXACTLY ONE)}

\textbf{UNRESOLVED (BUT STRICTLY ADVANCED).}

Round~9 gives a genuine strengthening beyond Round~8. We proved that the one-sided blow-up sets
\[
B_-:=\{x:\overline D_-G(x)=+\infty\},
\qquad
B_+:=\{x:\overline D_+G(x)=+\infty\}
\]
are not only residual (Round~7) but also $\mu$-full, where $\mu=dG$ is the singular distribution measure.
Hence
\[
\Xi:=\bigl(B_-\cap B_+\bigr)\setminus\mathcal E
\]
is a set that is simultaneously dense $G_\delta$, $\mu$-full, and disjoint from the special value set $\mathcal E$, and at every point of $\Xi$ both one-sided upper Dini derivatives are infinite.

So any hypothetical point where $G'(x)$ exists and is a finite positive number must now lie outside a set that is at once topologically generic, conull for the intrinsic distribution measure, and disjoint from the dense countable value set $\mathcal E$.
This is a strict advance, but it still stops short of a full proof or disproof of Erd\H{o}s's original question.

\subsection{8) COMPLETION ESTIMATE (MANDATORY)}

\textbf{COMPLETION: 93\%}

\subsection{9) REFERENCES}

\begin{enumerate}
\item I. J. Schoenberg, \emph{On asymptotic distributions of arithmetical functions}, Trans. Amer. Math. Soc. \textbf{39} (1936), 315--330.
\item P. Erd\H{o}s, \emph{On the distribution function of additive functions}, Ann. of Math. (2) \textbf{47} (1946), 1--20.
\item G. Tenenbaum and V. Toulmonde, \emph{Sur le comportement local de la r\'epartition de l'indicatrice d'Euler}, Funct. Approx. Comment. Math. \textbf{35} (2006), 321--338.
\item V. Toulmonde, \emph{Module de continuit\'e de la fonction de r\'epartition de $\varphi(n)/n$}, Acta Arith. \textbf{121} (2006), no.~4, 367--402.
\item P. Mattila, \emph{Geometry of Sets and Measures in Euclidean Spaces}, Cambridge Studies in Advanced Mathematics 44, Cambridge Univ. Press, 1995. (Used for the one-dimensional Besicovitch covering theorem.)
\end{enumerate}
